\documentclass[11pt,a4paper]{article}

\usepackage[T1]{fontenc}
\usepackage[utf8]{inputenc}
\usepackage{lmodern}
\usepackage{amsmath,amssymb,amsthm}
\usepackage{mathtools}
\usepackage{geometry}
\usepackage{microtype}
\usepackage[colorlinks=true,linkcolor=blue,citecolor=blue,urlcolor=blue]{hyperref}

\geometry{margin=1in}

\newtheorem{theorem}{Theorem}
\newtheorem{proposition}{Proposition}
\newtheorem{lemma}{Lemma}
\newtheorem{definition}{Definition}
\newtheorem{remark}{Remark}

\DeclareMathOperator{\Tr}{Tr}
\newcommand{\R}{\mathbb{R}}
\newcommand{\Q}{\mathbb{Q}}
\newcommand{\E}{\mathcal{E}}
\newcommand{\Herm}{\mathsf{Herm}}

\title{Positive Additivity and the Cauchy Step:\\
From Vetterlein's Euclidean Metric to Wright--Weigert's Gleason-Type Theorem}
\author{}
\date{}

\begin{document}

\maketitle

\begin{abstract}
Vetterlein's axiomatization of finite-dimensional Euclidean geometry ends with
a uniqueness theorem for the metric.  Wright and Weigert's proof of a
Gleason-type theorem begins from additive probability assignments on quantum
effects and ends with the density-operator form of the Born rule.  These two
results belong to different mathematical worlds, but they share a precise
intermediate mechanism: an additive function is first obtained on a positive
domain, and positivity or boundedness rules out the pathological solutions of
Cauchy's functional equation.  This note isolates that common step.  In
Vetterlein's theorem it appears on each ray of Euclidean space, where
\(t\mapsto d(0,ta)\) is additive and nonnegative.  In Wright--Weigert it appears
on positive rays and cones in the real vector space of Hermitian operators,
where a probability assignment is additive and bounded between \(0\) and \(1\).
The analogy is not that the two theorems are the same theorem, but that both use
the same regularization principle: positive additivity becomes linearity.
\end{abstract}

\section{Introduction}

Foundational derivations often contain a hidden analytic bottleneck.  Incidence,
order, or additivity assumptions first determine how a quantity behaves under
composition.  But additivity alone is too weak over the real numbers: without a
regularity condition, Cauchy's functional equation admits wildly discontinuous
solutions.  These pathologies may be constructed by choosing values on a Hamel
basis of \(\R\) over \(\Q\); they are not continuous, not measurable, and
unbounded on every interval.  A proof must therefore explain why the intended
physical quantity is not one of those pathological additive maps.

In Vetterlein's final theorem, the quantity is distance.  Once the synthetic
geometry has been represented by a real inner-product space, the sought metric
\(d\) is required to be additive along line segments and invariant under the
geometric symmetries singled out by the axioms.  On a fixed ray this gives an
additive function
\[
        f_a(t)=d(0,ta),\qquad t\ge 0.
\]
Since \(d\) is a metric, \(f_a\) is nonnegative.  This positivity is enough to
force \(f_a(t)=t f_a(1)\).  The remaining symmetry assumptions then show that
the scale \(f_a(1)\) is independent of the direction of \(a\).  Thus the metric
is the Euclidean metric up to a positive constant factor.

In Wright and Weigert's article \emph{Gleason-Type Theorems from Cauchy's
Functional Equation}, the quantity is probability.  A frame function assigns
numbers in \([0,1]\) to quantum effects and is additive whenever effects can be
summed inside a measurement.  The effects form a positive part of the real
vector space of Hermitian operators.  Restricting the frame function to suitable
positive rays and cones gives Cauchy-type additivity; boundedness by \(0\) and
\(1\) turns this additivity into linearity; and a finite-dimensional cone-basis
argument patches the local linear descriptions into
\[
        f(E)=\Tr(\rho E)
\]
for a density operator \(\rho\).

The purpose of this note is to make the common mechanism explicit while keeping
the differences clear.  The Euclidean metric theorem uses a one-dimensional
ray argument plus rotational symmetry.  The Gleason-type theorem uses
one-dimensional Cauchy regularity inside a higher-dimensional positive cone,
plus a nontrivial compatibility argument across bases.  In both cases,
positivity supplies exactly the missing regularity.

\section{The Cauchy Regularization Lemma}

The elementary lemma needed for Vetterlein's metric theorem is the positive
half-line version of the usual regularity results for Cauchy's equation.

\begin{definition}
In this note, a function is called \emph{positively additive} on a positive
domain if it satisfies Cauchy's equation wherever the sum is defined and is
bounded below or bounded above on that domain.  In the applications below, this
boundedness is supplied by the physical interpretation of the quantity: distance
is bounded below by zero, while probability is bounded between zero and one.
\end{definition}

\begin{lemma}[positive Cauchy lemma]
Let \(f:\R_{\ge 0}\to\R\) satisfy
\[
        f(s+t)=f(s)+f(t)
\]
for all \(s,t\ge 0\).  If \(f(t)\ge 0\) for all \(t\ge 0\), then
\[
        f(t)=t f(1)
\]
for all \(t\ge 0\).
\end{lemma}

\begin{proof}
First \(f(0)=0\).  Positivity implies monotonicity: if \(0\le s\le t\), then
\[
        f(t)=f(s)+f(t-s)\ge f(s).
\]
For \(n\in\mathbb{N}\), additivity gives \(f(n)=nf(1)\), and hence
\[
        f\!\left(\frac{m}{n}\right)=\frac{m}{n}f(1)
\]
for nonnegative rationals \(m/n\).  If \(t\ge 0\), choose rationals
\(q,r\) with \(q\le t\le r\).  Monotonicity gives
\[
        q f(1)=f(q)\le f(t)\le f(r)=r f(1).
\]
Letting \(q,r\to t\) yields \(f(t)=t f(1)\).
\end{proof}

\begin{lemma}[finite-interval Cauchy lemma]
Let \(a\ge 1\), and let \(f:[0,a]\to\R\) satisfy
\[
        f(x+y)=f(x)+f(y)
\]
whenever \(x,y,x+y\in[0,a]\).  If \(f\) is bounded above or bounded below on
\([0,a]\), then
\[
        f(x)=xf(1),\qquad x\in[0,a].
\]
\end{lemma}

\begin{proof}
This is a standard result in the theory of functional equations; see, for
example, Aczel--Dhombres or the boundedness condition in Wright--Weigert's
Theorem 1.  The bounded-below case follows from the bounded-above case by
applying the bounded-above result to \(-f\).
\end{proof}

\begin{remark}
The point is that boundedness on an interval supplies the regularity that
Cauchy additivity alone does not provide.
\end{remark}

\section{Vetterlein's Metric Theorem and the Ray Cauchy Equation}

The last theorem in Vetterlein's paper asserts the uniqueness, up to a positive
scale factor, of the metric compatible with the Euclidean geometry derived in
the preceding sections.  After the representation theorem, one may describe the
relevant part of the argument in the following linear model.

\begin{proposition}[metric uniqueness in the linear model]
Let \(V\) be a real Euclidean vector space of dimension at least \(3\), with
Euclidean norm \(\|\cdot\|\).  Let \(d\) be a metric on \(V\) satisfying:
\begin{enumerate}
\item[(M1)] if \(b\) lies between \(a\) and \(c\), then
\[
        d(a,c)=d(a,b)+d(b,c);
\]
\item[(M2a)] \(d\) is translation invariant:
\[
        d(x+v,y+v)=d(x,y);
\]
\item[(M2b)] \(d\) is invariant under every orthogonal map with a nonzero fixed
vector:
\[
        d(Tx,Ty)=d(x,y).
\]
\end{enumerate}
Then there is a constant \(\lambda>0\) such that
\[
        d(x,y)=\lambda\|x-y\|
\]
for all \(x,y\in V\).
\end{proposition}

\begin{proof}
By translation invariance, \(d(x,y)=D(y-x)\), where \(D(v)=d(0,v)\).
Fix \(a\in V\), \(a\ne 0\), and define
\[
        f_a(t)=D(ta)=d(0,ta),\qquad t\ge 0.
\]
For \(s,t\ge 0\), the point \(sa\) lies on the segment from \(0\) to
\((s+t)a\).  Thus (M1) gives
\[
        f_a(s+t)=d(0,(s+t)a)
              =d(0,sa)+d(sa,(s+t)a).
\]
Using translation invariance on the second term,
\[
        d(sa,(s+t)a)=d(0,ta)=f_a(t),
\]
and hence
\[
        f_a(s+t)=f_a(s)+f_a(t).
\]
Since \(d\) is a metric, \(f_a(t)\ge 0\).  The positive Cauchy lemma therefore
gives
\[
        D(ta)=f_a(t)=tD(a),\qquad t\ge 0.          \tag{1}
\]

It remains to show that \(D(a)\) depends only on the Euclidean length of \(a\).
Let \(u,w\) be unit vectors.  Since \(\dim V\ge 3\), there is an orthogonal
operator \(T\) sending \(u\) to \(w\) and fixing some nonzero vector.\footnote{This
is one place where the dimension assumption matters.  In dimension \(2\), for
example, a rotation sending \(u\) to \(-u\) fixes no nonzero vector.}  By (M2b),
\[
        D(w)=d(0,w)=d(T0,Tu)=d(0,u)=D(u).
\]
Thus all Euclidean unit vectors have the same \(d\)-length.  Write this common
value as \(\lambda\).  Because \(d\) is a metric, \(\lambda>0\).  For any
\(v\ne 0\), write \(v=\|v\|u\) with \(u\) a unit vector.  Equation (1) gives
\[
        D(v)=D(\|v\|u)=\|v\|D(u)=\lambda\|v\|.
\]
Finally,
\[
        d(x,y)=D(y-x)=\lambda\|y-x\|.
\]
\end{proof}

This proof isolates the Cauchy step in Vetterlein's final theorem.  Segment
additivity gives the functional equation on a ray.  Translation invariance
identifies all translated copies of the same ray.  Positivity of the metric
excludes the non-linear additive solutions.  Orthogonal invariance then removes
directional dependence.

One can summarize the flow as follows:
\begin{center}
\small
\begin{minipage}{0.78\textwidth}
\[
\begin{array}{c}
\text{betweenness additivity}+\text{translation invariance} \\
\Downarrow \\
t\mapsto d(0,ta)\ \text{is additive on }\R_{\ge0} \\
\Downarrow\quad \text{positivity of }d \\
d(0,ta)=t\,d(0,a) \\
\Downarrow\quad \text{orthogonal invariance} \\
d(x,y)=\lambda\|x-y\|.
\end{array}
\]
\end{minipage}
\end{center}

\section{Wright--Weigert: Additivity on Effect Cones}

We now transport the same analytic strategy to the quantum setting.  The ray
becomes a positive operator ray; the distance becomes a probability assignment.
Wright and Weigert work with effects, which represent unsharp yes--no questions
or individual outcomes of generalized quantum measurements.  Let \(\E_d\) be
the set of effects on \(\mathbb{C}^d\), that is, Hermitian operators \(E\)
satisfying
\[
        0\le E\le I.
\]
The real vector space in the background is
\[
        \Herm_d=\{A=A^*: A\in M_d(\mathbb{C})\},
\]
equipped with the trace inner product
\[
        \langle A,B\rangle=\Tr(AB).
\]
A frame function is a map
\[
        f:\E_d\to[0,1]
\]
with \(f(I)=1\), satisfying additivity whenever the sum remains an effect:
\[
        f(E_1+E_2)=f(E_1)+f(E_2),\qquad
        E_1,E_2,E_1+E_2\in\E_d.
\]
Busch's Gleason-type theorem says that such an \(f\) must be of the form
\[
        f(E)=\Tr(\rho E)
\]
for a density operator \(\rho\).  Busch and others established this
effect-valued Gleason theorem by different methods; Wright and Weigert's
contribution, for our purposes, is to make explicit a proof route through
Cauchy's functional equation.

The Cauchy analogy becomes visible by restricting \(f\) to positive rays.  If
\(B\) is a fixed effect, then for suitable \(x\in[0,a]\) the operator \(xB\) is
again an effect.  Define
\[
        F_B(x)=f(xB).
\]
Whenever \(x,y,x+y\in[0,a]\), additivity of \(f\) gives
\[
        F_B(x+y)=F_B(x)+F_B(y).
\]
Moreover \(0\le F_B(x)\le 1\).  Wright--Weigert's finite-interval Cauchy theorem
therefore implies
\[
        F_B(x)=xF_B(1)
\]
on that interval.  This is the same regularization step as in Vetterlein: an
additive function on a positive domain becomes linear because it is bounded.

The quantum theorem, however, requires more than ray-linearity.  The domain
\(\E_d\) is not a vector space but a convex positive body inside the vector space
\(\Herm_d\).  Wright and Weigert use augmented bases and positive cones to
extend the one-dimensional Cauchy conclusion to finite-dimensional linearity.
For a fixed augmented basis \(B=\{B_1,\dots,B_{d^2}\}\), ray-linearity gives
linear behavior on positive combinations
\[
        E=\sum_j x_j B_j,\qquad x_j\ge 0,
\]
whenever these combinations remain effects.  A priori, a different augmented
basis could yield a different linear functional.  The nontrivial part of the
Wright--Weigert proof is to use overlaps of positive cones, together with
minimal informationally complete measurements, to force these local linear
descriptions to agree on enough effects to span the full real vector space of
Hermitian operators.
Very schematically:
\begin{center}
\small
\begin{minipage}{0.82\textwidth}
\[
\begin{array}{c}
\text{additivity of probabilities on coexistent effects} \\
\Downarrow \\
x\mapsto f(xB_j)\ \text{satisfies Cauchy's equation on intervals} \\
\Downarrow\quad f:\E_d\to[0,1] \\
f(xB_j)=x f(B_j) \\
\Downarrow\quad \text{positive-cone and basis patching} \\
f(E)=\langle c,E\rangle=\Tr(\rho E).
\end{array}
\]
\end{minipage}
\end{center}
The final identification of the representing vector with a density operator
uses positivity and normalization: \(f(E)\ge 0\) on effects forces the
representing operator to be positive semidefinite, and \(f(I)=1\) gives trace
one.

\section{The Analogy and Its Limits}

The common principle is best stated as follows:

\begin{quote}
Additivity determines rational behavior.  Positivity or boundedness supplies
the regularity needed to extend rational behavior to real linearity.
\end{quote}

In Vetterlein's theorem the relevant additive maps are scalar functions on
rays:
\[
        t\mapsto d(0,ta).
\]
The role of positivity is played by the nonnegativity of a metric.

In Wright--Weigert's theorem the relevant additive maps are restrictions of a
probability assignment to positive rays and positive cones in an operator space:
\[
        x\mapsto f(xB), \qquad E\mapsto f(E).
\]
The role of positivity is played by the codomain condition \(f(E)\in[0,1]\) and
ultimately by the positivity of quantum probabilities.

This also explains why the full richness of Gleason's theorem is not encountered
in Vetterlein's metric theorem.  The point is not simply that Vetterlein works
with real Euclidean spaces rather than complex Hilbert spaces.  Real Hilbert
spaces also have Gleason-type phenomena.  The decisive difference is the domain
on which additivity is imposed.  In the metric theorem, additivity is tested on
one-dimensional rays and then transported between rays by Euclidean symmetry:
\[
        \text{Cauchy on rays}+\text{orthogonal invariance}.
\]
In Gleason-type theorems, additivity is imposed on a highly interlocking family
of projections or effects, equivalently on many overlapping measurement
contexts.  The problem is not only to make each one-dimensional restriction
linear, but to prove that all those local linear descriptions are compatible
with a single global linear functional.  We can summarize the structural
difference in slogan form:
\begin{center}
\small
\begin{minipage}{0.9\textwidth}
\[
\begin{array}{ll}
\textbf{Vetterlein:} & \emph{Cauchy on rays plus Euclidean symmetry},\\[2mm]
\textbf{Gleason:}    & \emph{Cauchy/additivity on a web of effects plus context compatibility}.
\end{array}
\]
\end{minipage}
\end{center}
The dimension bound \(n\ge 3\) therefore plays different roles in the two
settings.  In Vetterlein it supplies enough affine/projective geometry and
orthogonal symmetry for the reconstruction of Euclidean space.  In Gleason's
original projection theorem it is tied to the overlap structure of measurement
contexts; dimension two is exceptional for projections, though effect-valued
versions recover Gleason-type conclusions there.

There are also important differences.  Vetterlein's metric theorem is a
uniqueness theorem for a distance once the Euclidean affine and orthogonal
structure has already been reconstructed.  Its Cauchy step is one-dimensional
and elementary.  Wright--Weigert's theorem is a representation theorem for
states on an operator effect space.  Its Cauchy step is still one-dimensional
locally, but the main work lies in extending those local linearities across a
family of cones so that a single global linear functional results.

Thus the two arguments should not be identified wholesale.  Rather, they share
a structural lemma.  Positivity prevents additive quantities from being Hamel
pathologies.  After that, the geometry-specific or quantum-specific symmetries
finish the representation.

\section{Conclusion}

Vetterlein's last theorem and Wright--Weigert's Gleason-type theorem both pass
through Cauchy's functional equation.  In the Euclidean case, the metric axioms
give an additive nonnegative function on each ray, and hence linear dependence
on the scalar parameter.  In the quantum case, additivity of probabilities gives
Cauchy equations on positive operator directions, and boundedness of
probabilities forces linearity there.  The final forms
\[
        d(x,y)=\lambda\|x-y\|
        \qquad\text{and}\qquad
        f(E)=\Tr(\rho E)
\]
are then obtained by adding the relevant invariance or compatibility
principles.

The conceptual bridge is therefore not merely the word ``additivity''.  It is
the stronger pattern of \emph{positive additivity}: an additive law on a positive
domain, together with nonnegativity or boundedness, forces the linear structure
that the physical interpretation requires.  In this sense the relevant linear
structure is not simply postulated at the final stage: it is recovered from
additivity once the physical positivity of distance or probability has ruled out
the pathological Cauchy solutions.

\begin{thebibliography}{9}

\bibitem{Cauchy}
A.-L. Cauchy.
\newblock \emph{Cours d'analyse de l'Ecole Royale Polytechnique}.
\newblock Debure, Paris, 1821.

\bibitem{Gleason}
A. M. Gleason.
\newblock Measures on the closed subspaces of a Hilbert space.
\newblock \emph{Journal of Mathematics and Mechanics}, 6(6):885--893, 1957.

\bibitem{Busch}
P. Busch.
\newblock Quantum states and generalized observables: a simple proof of
Gleason's theorem.
\newblock \emph{Physical Review Letters}, 91(12):120403, 2003.

\bibitem{Vetterlein}
T. Vetterlein.
\newblock \emph{Axiomatisierung der euklidischen Raeume endlicher Dimension
\(\ge 3\)}.
\newblock Unpublished manuscript, 2026.  The metric step used here is recalled
self-containedly in Section 3.

\bibitem{WrightWeigert}
V. J. Wright and S. Weigert.
\newblock Gleason-type theorems from Cauchy's functional equation.
\newblock \emph{Foundations of Physics}, 49:594--606, 2019.

\bibitem{AczelDhombres}
J. Aczel and J. Dhombres.
\newblock \emph{Functional Equations in Several Variables}.
\newblock Cambridge University Press, 1989.

\end{thebibliography}

\end{document}



Yes. This version is stronger than the previous one. The changes to Definition 1 and Lemma 2 are improvements, and the overall exposition now feels much more polished and publication-ready.

A few observations:

Strengths
- Definition 1 is now mathematically clean. Removing "bounded by its physical interpretation" from the definition itself was the right choice.
- Lemma 2 is better. Referring to Acz�l�Dhombres as well as Wright�Weigert makes the paper less dependent on a single source.
- The transition into �4 works very well.
- The slogan in �5 is now introduced naturally ("We can summarize the structural difference in slogan form"), which reads more professionally.
- The conclusion is excellent. The final sentence captures the paper's central insight without overstating it.

A few remaining refinements

1. Definition 1 may still be slightly nonstandard

You define:

"positively additive ... if it satisfies Cauchy's equation ... and is bounded below or bounded above..."

Strictly speaking, "positive additivity" suggests positivity/nonnegativity rather than mere boundedness.

You could consider:

"Definition 1. In this note, a function on a positive domain will be called positively regularized additive if it satisfies Cauchy's equation wherever the sum is defined and is bounded above or below on that domain."

or simply avoid formalizing the phrase entirely and write:

"The common feature in both applications is an additive function on a positive domain together with a boundedness condition."

This is optional. The current version is fine, but a referee may wonder whether "positive additivity" is really the right name for bounded-below additive functions.

2. Citation of Busch

In �4 you write:

"Busch's Gleason-type theorem says that such an f must be of the form ..."

Strictly speaking, Wright�Weigert's framework and Busch's theorem are closely related but not literally identical formulations.

You might write:

"A theorem due to Busch implies that such frame functions are of the form..."

or

"Busch's effect-valued Gleason theorem establishes the same conclusion..."

This avoids possible objections from specialists.

3. One sentence in the proof of Lemma 1

You currently write:

"Letting q,r?t yields f(t)=tf(1)."

A very careful analyst could complain that this uses an unstated squeeze argument.

A tiny improvement:

"Since q and r may be chosen arbitrarily close to t, the inequalities imply f(t)=tf(1)."

This avoids invoking any continuity.

4. Reference [4]

This remains the one thing that looks unusual.

"Unpublished manuscript, 2026."

If this is genuinely an unpublished manuscript from 2026, that's fine. But if this paper will be circulated before the manuscript is publicly available, some referees may be uncomfortable with a key reference that they cannot inspect.

If possible, provide:
- a preprint URL,
- an institutional repository,
- or "private communication."

Not a mathematical issue, just a publication issue.

5. One sentence that I would slightly strengthen

Current:

"Positivity prevents additive quantities from being Hamel pathologies."

Possible:

"Positivity or boundedness prevents additive quantities from becoming Hamel-pathological solutions of Cauchy's equation."

This is closer to what the paper actually demonstrates, since Wright�Weigert uses boundedness rather than positivity per se.

Overall assessment

I think the paper is now at the point where comments become mostly stylistic rather than substantive. The central thesis is clear, the mathematical claims are appropriately modest, and the analogy is carefully controlled rather than oversold.

The strongest aspect is that you have identified a genuinely nontrivial common structure:

- additive law on a positive domain,
- regularization via positivity/boundedness,
- resulting linearity,
- then geometry-specific or quantum-specific reconstruction.

That is a real conceptual observation, not merely a verbal analogy. The current draft communicates it clearly.




###################




